\section*{Bab \ref{ch:b3:dna-repair} --- Kemantapan Genom: Kerusakan, Perbaikan, dan Rekombinasi DNA}

\begin{solution}{exo:b3:dna-repair:1}
Urasil: spontan (\omterm{def:b3:dna-repair:lesions}{deaminasi} C), \omterm{def:b3:dna-repair:excision}{perbaikan eksisi basa} oleh
\omterm{def:b3:dna-repair:excision}{glikosilase} urasil-DNA. Dimer timina: lingkungan (UV), \omterm{def:b3:dna-repair:excision}{perbaikan eksisi
nukleotida}. \omterm{def:b3:dna-repair:lesions}{8-oksoguanina}: spontan (oksidasi), \omterm{def:b3:dna-repair:excision}{perbaikan eksisi
basa} oleh \omterm{def:b3:dna-repair:excision}{glikosilase} 8-oksoG. Salah pasang G--T di belakang garpu:
galat replikasi, \omterm{prop:b3:dna-repair:fidelity}{perbaikan salah pasang}. \omterm{def:b3:dna-repair:lesions}{Putus untai ganda} pada G1:
lingkungan (radiasi) atau spontan; \omterm{def:b3:dna-repair:dsb}{penyambungan ujung tak homolog},
karena tidak ada kromatid saudari yang tersedia.
\end{solution}

\begin{solution}{exo:b3:dna-repair:2}
Eksisi nukleotida mengenali sifat fisik yang dimiliki bersama oleh semua
\omterm{def:b3:dna-repair:lesions}{lesi} ruah --- kerusakan bentuk heliksnya, atau RNA polimerase yang terhenti
--- lalu membuang tambalan yang memuat apa pun penyebabnya. \omterm{def:b3:dna-repair:lesions}{Lesi} kecil
(urasil, 8-oksoG, basa teralkilasi) nyaris tidak merusak bentuk heliks,
sehingga semuanya harus dikenali dari kimianya, satu \omterm{def:b3:dna-repair:excision}{glikosilase} tiap macam
basa yang salah.
\end{solution}

\begin{solution}{exo:b3:dna-repair:3}
Sebuah salah pasang terdiri atas dua basa normal, dan hanya basa pada untai
yang baru yang keliru; sistemnya harus tahu untai mana yang baru.
\emph{E.~coli} menandai untai tetuanya dengan gugus metil pada GATC,
yang belum dimiliki untai barunya selama beberapa menit; MutH menakik untai
yang tak termetilasi. Pada mutan \emph{dam} tidak ada untai yang
termetilasi: MutH menakik untai mana saja, perbaikannya membuang basa yang
benar sesering basa yang salah --- sehingga separuh galatnya justru dipatri
menjadi mutasi (sebuah mutator) --- dan takik pada kedua untai di tapak GATC
yang berdekatan menimbulkan \omterm{def:b3:dna-repair:lesions}{putus untai ganda}.
\end{solution}

\begin{solution}{exo:b3:dna-repair:4}
Kira-kira $35\times 2 = 70$ \omterm{def:b3:dna-repair:lesions}{putus untai ganda}. Sesudah sejam sebagian besar
telah tersambung: dengan waktu paruh mendekati \qty{30}{min}, sekitar
$70/4 \approx 18$ fokus. Sesudah enam jam hampir semuanya telah diperbaiki
($70/2^{12} < 1$), sehingga hanya tersisa beberapa putus rumit yang bertahan.
\end{solution}

\begin{solution}{exo:b3:dna-repair:5}
(a) $10^{-5}\times 10^{-2}\times 10^{-3} = 10^{-10}$; $6.4\times
10^{9}\times 10^{-10} = 0.64$ mutasi tiap pembelahan. (b) Tanpa
\omterm{prop:b3:dna-repair:fidelity}{perbaikan salah pasang} $10^{-7}$; $640$ tiap pembelahan. (c) Tanpa \omterm{prop:b3:dna-repair:fidelity}{uji baca}
$10^{-8}$; $64$ tiap pembelahan. Mutan \omterm{prop:b3:dna-repair:fidelity}{perbaikan salah pasang} itulah yang
lebih berbahaya, seribu kali lipat di atas normal.
\end{solution}

\begin{solution}{exo:b3:dna-repair:6}
$k = \ln 2/\qty{0.25}{h} = \qty{2.8}{h^{-1}}$; $D = 10^{4}/24 =
\qty{420}{h^{-1}}$; $L^{*} = 420/2.8 \approx 150$ tapak tanpa basa hadir
pada satu saat. Perbaikan sepuluh kali lebih lambat: $L^{*} \approx \num{1500}$.
Garpunya melewati tiap lokus sekali, dan berjumpa dengan \omterm{def:b3:dna-repair:lesions}{lesi} yang hadir di
situ pada saat itu; dirata-ratakan atas genomnya ia berjumpa dengan sekitar
$L^{*}$ \omterm{def:b3:dna-repair:lesions}{lesi} --- sekitar $150$ pada keadaan normal, $\num{1500}$ dengan obatnya.
\end{solution}

\begin{solution}{exo:b3:dna-repair:7}
$k_{\text{NHEJ}} = \ln 2/0.5 = \qty{1.39}{h^{-1}}$, $k_{\text{HR}} =
\ln 2/4 = \qty{0.17}{h^{-1}}$; bagian HR-nya $0.17/1.56 = 0.11$;
waktu paruhnya $\ln 2/1.56 = \qty{27}{min}$. Garpu yang runtuh memberi
putus \emph{berujung satu}: tidak ada ujung kedua untuk disambungkan NHEJ,
sehingga hanya rekombinasi dengan saudarinya yang dapat memulihkan garpu itu
--- lagi pula ia satu-satunya jalur yang saksama.
\end{solution}

\begin{solution}{exo:b3:dna-repair:8}
(a) Lynch: mikrosatelitnya tak mantap (panjangnya berbeda antarsel
tumor), kanker kolorektal, endometrium, lambung, dan ovarium, tanggapan
normal terhadap cahaya matahari. (b) XP-A: mikrosatelitnya mantap, lepuh
matahari hebat pada paparan minimal, kanker kulit dan mata pada masa
kanak-kanak, dan pada golongan ini kemunduran saraf yang berangsur. (c) XP-V:
mikrosatelitnya mantap, perbaikan eksisinya normal sehingga lepuh mataharinya
lebih ringan, tetapi dimer yang lolos dari eksisi dilangkahi polimerase rawan
galat --- kanker kulit, yang muncul lebih lambat daripada pada golongan~A.
\end{solution}

\begin{solution}{exo:b3:dna-repair:9}
Sel hati: Cre telah mengeksisi ekson~3 dari kedua alelnya (sebagai
lingkaran yang hilang), sehingga gennya tak aktif di setiap hepatosit sejak
lahir --- penidakaktifan yang khas hati. Neuron: tanpa Cre, alel ber-loxP-nya
utuh dan berfungsi. Arah berlawanan: Cre membalikkan eksonnya lalu
membalikkannya kembali tanpa henti, sehingga pada saat mana pun kira-kira
separuh selnya mengusung ekson terbalik yang tak berfungsi --- sebuah mozaik,
bukan penghapusan yang bersih.
\end{solution}

\begin{solution}{exo:b3:dna-repair:10}
Seiring LexA dipotong, kepekatannya turun berangsur. Operator yang
mengikat LexA dengan lemah dikosongkan pada penurunan kecil --- \emph{uvrA},
\emph{uvrB}, \emph{recA} --- sehingga perbaikan yang saksama mulai lebih
dahulu; operator \emph{sulA} dan operator polimerase rawan galat mengikat
LexA dengan kuat dan baru dikosongkan ketika LexA nyaris habis, yakni ketika
kerusakannya telah bertahan lama. Bakteri tanpa polimerase rawan
galat akan mati pada garpu yang terhalang \omterm{def:b3:dna-repair:lesions}{lesi} yang tak dapat
dilangkahinya, dan, tanpa mutagenesis imbasan, akan mengembangkan kekebalan
terhadap antibiotiknya jauh lebih lambat --- merugikan bagi bakterinya,
menguntungkan bagi pasiennya; penghambat \omterm{prop:b3:dna-repair:sos}{tanggapan SOS} dikaji karena alasan
itu.
\end{solution}

\begin{solution}{exo:b3:dna-repair:11}
Rantainya: PARP dihambat $\to$ putus untai tunggal bertahan $\to$ sebuah
garpu replikasi mengubah masing-masingnya menjadi \omterm{def:b3:dna-repair:lesions}{putus untai ganda} berujung satu
$\to$ perbaikannya menuntut \omterm{def:b3:dna-repair:dsb}{rekombinasi homolog} $\to$ sel tumornya, yang
kekurangan \omterm{def:b3:dna-repair:dsb}{BRCA}, tidak sanggup $\to$ salah sambung, hilangnya kromosom,
kematian. Sel normal, dengan satu alel yang berfungsi, berekombinasi dan
bertahan. Kekebalan: (1) mutasi kedua pada \emph{\omterm{def:b3:dna-repair:dsb}{BRCA}} yang memulihkan
kerangka baca dan protein yang bekerja --- dideteksi dengan mengurutkan
tumornya atau DNA tumor yang beredar di dalam darah, dan dengan kembalinya
fokus \omterm{def:b3:dna-repair:dsb}{Rad51} sesudah penyinaran; (2) hilangnya protein yang menghalangi
pemangkasan ujung (53BP1--shieldin), yang sebagian memulihkan rekombinasi
tanpa BRCA1 --- dideteksi dari ketiadaannya pada pewarnaan dan dari pulihnya
fokus \omterm{def:b3:dna-repair:dsb}{Rad51}; juga mutasi PARP1 yang mencegah terperangkapnya pada DNA, atau
pompa penyembur obat, yang berturut-turut dideteksi lewat pengurutan dan
ekspresi.
\end{solution}

\begin{solution}{exo:b3:dna-repair:12}
Dua duplek homolog, yang satu mengusung $A$ dan yang lain $a$, saling
mempertukarkan untai tunggal melintasi penandanya: masing-masing kini memiliki
satu untai dari tiap tetua sepanjang bentangan itu --- heterodupleks dengan
salah pasang $A/a$. Bila \omterm{prop:b3:dna-repair:fidelity}{perbaikan salah pasang} mengubah salah pasang pada
duplek yang diserbu menjadi $A$, tiga kromatid mengusung $A$ dan satu $a$:
$3{:}1$; bila diubah menjadi $a$: $1{:}3$; bila dibiarkan tanpa perbaikan,
kromatidnya memilah $A$ dan $a$ pada mitosis pertama sesudah meiosis, sehingga
memberi dua spora anak yang berbeda (pemilahan pascameiosis, $5{:}3$ pada
askus berspora delapan). Bentangan heterodupleksnya hanya beberapa ratus
sampai beberapa ribu pasangan basa di sekeliling putus yang memulainya,
sehingga hanya penanda di dalamnya yang dapat dikonversi, dan penanda itu,
menurut konstruksinya, berdampingan dengan pindah silangnya.
\end{solution}

\begin{solution}{pb:b3:dna-repair:1}
\textbf{1.} $10^{4} + 400 + 10^{4} \approx 2\times 10^{4}$ \omterm{def:b3:dna-repair:lesions}{lesi} tiap
hari, $7\times 10^{6}$ tiap tahun --- kira-kira satu tiap kilobasa genomnya
tiap tahun.
\textbf{2.} $10^{-5}\times 10^{-2}\times 10^{-3} = 10^{-10}$ tiap bp;
$6.4\times 10^{9}\times 10^{-10} = 0.64$ mutasi tiap pembelahan.
\textbf{3.} $0.64\times 0.015 \approx 0.01$ mutasi penyandi tiap
pembelahan; sepanjang $50$ pembelahan, sekitar $0.5$.
\textbf{4.} $10^{-7}$ tiap bp, $640$ mutasi tiap pembelahan. Sel
tergelincir: tanpa perbaikan $10^{-4}\times 40\times 10^{9} = 4\times 10^{6}$;
dengan perbaikan $10^{-7}\times 40\times 10^{9} = 4000$ --- selisih seribu
kali lipat, yang kasatmata sebagai \omterm{def:b3:dna-repair:mmr}{ketakmantapan mikrosatelit}.
\textbf{5.} Urasil asing bagi DNA, dikenali \omterm{def:b3:dna-repair:excision}{glikosilase} urasil-DNA
lalu dieksisi: diperbaiki. Timina dari 5-metilsitosina adalah
basa normal yang berhadapan dengan G; tak ada yang menandainya salah di luar
replikasi, dan ia menjadi mutasi C$\to$T --- mekanisme yang telah menguras
\omterm{def:b3:chromatin-epigenetics:methylation}{CpG} dari genomnya.
\textbf{6.} \omterm{def:b3:dna-repair:excision}{Glikosilase} 8-oksoG membuang basa teroksidasi itu selagi ia
masih berhadapan dengan C (sebelum replikasi); \omterm{def:b3:dna-repair:excision}{glikosilase} kedua membuang
A yang salah disisipkan berhadapan dengan 8-oksoG (sesudah replikasi, sehingga
memulihkan peluang memperbaiki G-nya); hidrolasenya memusnahkan dGTP
teroksidasi sehingga tidak disisipkan berhadapan dengan A. Ketika ketiganya
gagal, 8-oksoG berpasangan dengan A dan putaran berikutnya mematri
transversi G$\cdot$C$\to$T$\cdot$A.
\textbf{7.} Normal $k = \ln 2/2 = \qty{0.35}{h^{-1}}$; XP $k = \ln 2/70
= \qty{0.0099}{h^{-1}}$.
\textbf{8.} Normal $10^{5} e^{-2.77} \approx 6200$; XP $10^{5}
e^{-0.079} \approx \num{92000}$.
\textbf{9.} Mutasinya $6.2$ dan $92$ tiap sel; penyandinya $0.09$ dan $1.4$.
\textbf{10.} $500$ paparan: $47$ mutasi penyandi (normal) dan $690$
(XP), lawan $0.5$ dari replikasi --- cahaya matahari mendominasi
muatan mutasi kulit bahkan pada orang normal.
\textbf{11.} Urutan penyandinya $9.6\times 10^{7}$ bp; gen pemacunya
sebesar $4\times 10^{4}/9.6\times 10^{7} = 4.2\times 10^{-4}$ dari itu.
Rata-rata terkena $690\times 4.2\times 10^{-4} = 0.29$; $P(\ge 1) = 1 - e^{-0.29} =
0.25$. Dari $10^{9}$ sel, $2.5\times 10^{8}$ mengusung mutasi pemacu
(anak normal: rata-rata $0.02$, \qty{2}{\%}).
\textbf{12.} Sebuah dimer bukan mutasi; ia baru menjadi mutasi ketika
sebuah garpu menyalinnya lewat pelangkahan yang rawan galat, sehingga sel
yang membelah dengan banyak dimer tak terperbaiki itulah yang bermutasi.
Cahaya ultraungu hanya mencapai kulit dan mata; jaringan dalamnya tidak
menerima dimer dan, pada XP, hampir normal.
\textbf{13.} $70$ putus; sesudah \qty{1}{h} dengan NHEJ saja, $70\times
2^{-2} \approx 18$ fokus.
\textbf{14.} $k_{\text{NHEJ}} = \qty{1.39}{h^{-1}}$, $k_{\text{HR}} =
\qty{0.17}{h^{-1}}$; bagian HR-nya $0.17/1.56 = 0.11$.
\textbf{15.} Waktu paruhnya $\ln 2/1.56 = \qty{27}{min}$; sesudah \qty{1}{h}
tersisa $70\,e^{-1.56} \approx 15$ putus.
\textbf{16.} $70\times 0.015\times 0.7 \approx 0.7$ gen dirusak tiap
sel.
\textbf{17.} $n = 70$: $2415$ pasangan $\times 3\times 10^{-4} = 0.72$
translokasi. $n = 7$: $21\times 3\times 10^{-4} = 0.006$. Cacah putusnya
$\propto D$ tetapi cacah \emph{pasangan} putus yang serentak
$\propto n^{2} \propto D^{2}$.
\textbf{18.} $D = \qty{7}{h^{-1}}$, $L^{*} = 7/1.39 \approx 5$ putus
hadir pada satu saat: $10$ pasangan, $0.003$ translokasi --- dua ratus
kali lebih sedikit daripada dosis yang sama sekaligus. Pemenggalan dosis
memberi jaringan normal kesempatan memperbaiki di antara penggalnya dan
menjaga putus yang serentak tetap sedikit.
\textbf{19.} Perbaikan utuh: $e^{-2/1.5} = 0.26$ pada \qty{2}{Gy},
$e^{-4} = 0.018$ pada \qty{6}{Gy}. Tanpa NHEJ: $e^{-4} = 0.018$ dan
$e^{-12} = 6\times 10^{-6}$. $D_{0}$ adalah dosis yang rata-rata
meninggalkan satu \omterm{def:b3:dna-repair:lesions}{lesi} mematikan tak terperbaiki tiap sel (ketahanan $1/e$).
\textbf{20.} \qty{11}{\%} putus G2 berujung dua yang tadinya akan
diperbaiki HR secara saksama kini lewat NHEJ, dan setiap putus garpu
berujung satu --- yang sama sekali tak dapat diperbaiki NHEJ secara benar ---
salah disambung: sepanjang bertahun-tahun genomnya menimbun delesi,
translokasi, dan perubahan cacah salinan, yakni parut khas tumor \emph{\omterm{def:b3:dna-repair:dsb}{BRCA}}.
\textbf{21.} $10^{4}\times 0.01 = 100$ putus berujung satu tiap hari.
Putus berujung satu tidak memiliki ujung pasangan; NHEJ hanya dapat
membiarkannya atau menyambungkannya ke ujung lain yang tak berkaitan,
sehingga menghasilkan translokasi atau kehilangan.
\textbf{22.} Satu alel \emph{BRCA2} yang berfungsi membuat protein yang cukup
bagi rekombinasi: sel normalnya memperbaiki putus garpunya dan bertahan
terhadap penghambatnya.
\textbf{23.} Pergeseran kerangka kedua di hilir yang pertama memulihkan
kerangka bacanya dan menghasilkan BRCA2 yang lebih pendek tetapi sebagian
berfungsi, sehingga rekombinasinya kembali dan obatnya tak lagi membunuh.
Dideteksi dengan mengurutkan \emph{BRCA2} pada tumor yang kebal atau pada DNA
tumor yang beredar, dan secara fungsional dari munculnya kembali fokus \omterm{def:b3:dna-repair:dsb}{Rad51}
sesudah penyinaran.
\textbf{24.} Pembunuhannya menuntut \omterm{prop:b3:dna-repair:fidelity}{perbaikan salah pasang} untuk menggarap
basa yang teralkilasi; tumor Lynch yang tak memilikinya bersifat tahan ---
obatnya gagal membunuhnya (tumor semacam itu justru diobati dengan cara lain,
termasuk imunoterapi, yang muatan mutasinya yang berat justru membuatnya
peka).
\textbf{25.} Sekitar \num{92000} dimer pada garpu sel XP;
$10^{-10}$ mutasi tiap pasangan basa tiap pembelahan; HR memperbaiki
\qty{11}{\%} putus G2.
\end{solution}
